% !TEX root = ../main.tex

\section{\bilingual{Lemmas on Spine Random Walk}{有关脊柱随机游走的引理}}\label{section:spine-random-walk-lemmas}

\bilingual{
    The results from this section onward are based on the works of A\"id\'ekon and Hu \mycite{aïdékon2025averageweightpercolationcompletegraph}. Consider the random walk $(Y_n)$ obtained from \nameref{lemma:many-to-one}.
}{
    从本节开始的结论基于并推广了A\"id\'ekon与Hu \mycite{aïdékon2025averageweightpercolationcompletegraph} 的工作.考虑由\nameref{lemma:many-to-one}得到的随机游走$(Y_n)$.
}

\begin{lemma}[\bilingual{Small Deviation}{水平管道意义下的小偏差}, A\"id\'ekon--Hu \mycite{aïdékon2025averageweightpercolationcompletegraph}]\label{lemma:4.1}
    \bilingual{
        Let $v>0$ and $0<d<\frac{\pi^2}2\sigma^2$. There exists $\Delta_0>0$ such that for all $\Delta\ge\Delta_0$, $r_0\le 0\le r_0+\Delta$ and $L\ge 1$,
    }{
        设 $v>0$ 且 $0<d<\frac{\pi^2}2\sigma^2$. ~存在 $\Delta_0>0$ 使得对所有 $\Delta\ge\Delta_0$,~ $r_0\le 0\le r_0+\Delta$ 且 $L\ge 1$,
    }
    {
        \setlength{\belowdisplayskip}{0pt}
        \begin{equation}
            \mathbb P(r_0\le Y_j\le r_0+\Delta,\,\forall j\in[\![1,L]\!])\le\exp\left( \Delta^v-\frac{d}{\Delta^2}L\right).
            \label{eq:4.1}
        \end{equation}
    }
\end{lemma}

% \vspace{-15pt}
\begin{figure}[H]
    \centering
    \begin{tikzpicture}[
        scale=1.5,
        >=Stealth,
        % 随机游走路径样式
        rw-path/.style={thick, blue!70!black, mark=*, mark size=1.5pt},
        % 管道填充样式
        tube-fill/.style={gray!10},
        % 标注样式
        label-font/.style={font=\small}
    ]

        % 1. 绘制背景管道填充 (r_0 到 r_0 + delta)
        \fill[tube-fill] (0, -1.5) rectangle (4, 0.5);

        % 2. 绘制辅助虚线
        \draw[dashed, black!60] (0, -2) -- (4.5, -2) node[right] {$-\Delta$};
        % \draw[dashed, black!60] (0, 0) -- (4.5, 0) node[right] {$y = 0$};
        \draw[dashed, black!60] (4, -2.5) -- (4, 1) node[above] at (4, 1) {$L$};

        % 3. 绘制管道实线 (r_0 到 r_0 + delta)
        \draw[thick, black] (0, -1.5) -- (4.5, -1.5) node[right] {$r_0$};
        \draw[thick, black] (0, 0.5) -- (4.5, 0.5) node[right] {$r_0 + \Delta$};

        % 4. 绘制坐标轴 (Y轴实线)
        \draw[->, thick] (0, -2.5) -- (0, 1.2) node[above] {$Y_n$};
        \draw[->, thick, dashed] (0, 0) -- (5, 0) node[right] {$n$};

        % 5. 绘制随机游走 (起点在 0,0，始终保持在 [-1.5, 0.5] 之间)
        % 选取了一些看起来随机的点
        \draw[rw-path] plot coordinates {
            (0, 0) 
            (0.8, -0.6) 
            (1.5, 0.2) 
            (2.2, -1.1) 
            (3.1, -0.4) 
            (4, -0.8)
        };

        % 6. 标注 Delta (使用大括号)
        \draw[decorate, decoration={brace, amplitude=5pt}] 
            (-0.1, -1.5) -- (-0.1, 0.5) node[midway, left=6pt] {$\Delta$};

        % 标注特定的点
        % \node[below left] at (0,0) {$0$};

    \end{tikzpicture}
    \caption{\bilingual{Illustration of the small deviation of the random walk \texorpdfstring{$Y_n$}{Yn} within the horizontal tube \texorpdfstring{$[r_0, r_0+\Delta]$}{[r_0, r_0+Delta]}}{随机游走 \texorpdfstring{$Y_n$}{Yn} 在水平管道 \texorpdfstring{$[r_0, r_0+\Delta]$}{[r_0, r_0+Delta]} 内的小偏差示意图}}
\end{figure}

\begin{proof}\label{proof:lemma:4.1}
    \bilingual{
        By applying \nameref{theorem:mogulskii} (see \Cref{lemma:mogulskii-cal} (ii))% and \Cref{example:mogulskii-L-1})
        , for all $a\in(0,1)$ and $\ell=\ell(\Delta)=\lfloor\frac 1d\Delta^{2+v}\rfloor$,
    }{
        通过应用 \nameref{theorem:mogulskii} (参见\Cref{lemma:mogulskii-cal} (ii))% 和\Cref{example:mogulskii-L-1})
        ,对所有 $a\in(0,1)$ 和 $\ell=\ell(\Delta)=\lfloor\frac 1d\Delta^{2+v}\rfloor$,~
    }
    \begin{equation*}
        \lim_{\Delta\to\infty}
        \frac{\Delta^2}{\ell}
        \log\mathbb P\left(
                -a\Delta\le Y_j\le(1-a)\Delta,\,
                \forall j\in[\![1,\ell]\!]
        \right)
        =
        -\frac{\pi^2}2\sigma^2.
        \label{eq:lemma:4.1:mogulskii-applied}
    \end{equation*}

    \bilingual{
        Fix $m\ge 1$ so that $d{(1+\frac 1m)}^2<\frac{\pi^2}2\sigma^2$. Apply the above equation to $(1+\frac 1m)\Delta$ in lieu of $\Delta$ and
    }{
        固定 $m\ge 1$ 使得 $d{(1+\frac 1m)}^2<\frac{\pi^2}2\sigma^2$. ~应用上式,应用时将上式中的 $\Delta$ 替换为 $(1+\frac 1m)\Delta$,~ 并且令
    }
    \begin{equation*}
        a=\frac{1-\frac km}{1+\frac 1m},\quad
        k\in[\![0,m-1]\!],
    \end{equation*}
    \bilingual{
        and we can get, for all $k\in[\![0,m-1]\!]$,
    }{
        从而得,对所有 $k\in[\![0,m-1]\!]$,~
    }
    \begin{align*}
        &\quad\lim_{\Delta\to\infty}
        \frac 
        {
            {\left(1+\frac 1m\right)}^2
            \Delta^2
        } \ell
        \log\mathbb P\left(
                \left(-1+\frac km\right)\Delta
                \le Y_j
                \le\frac{k+1}{m}\Delta,\,
                \forall j\in[\![1,\ell]\!]
        \right)
        \\ &=
        -\frac{\pi^2}{2}\sigma^2
        <-d{\left(1+\frac 1m\right)}^2.
    \end{align*}
    \bilingual{
        Then we find $\Delta_0=\Delta(d)>0$ such that for all $\Delta\ge\Delta_0$ and $k\in[\![0,m-1]\!]$,
    }{
        于是可得 $\Delta_0=\Delta(d)>0$ 使得对所有 $\Delta\ge\Delta_0$ 和 $k\in[\![0,m-1]\!]$,~ 
    }
    \begin{equation*}
        \mathbb P\left(
            \left(-1+\frac km\right)\Delta
            \le Y_j
            \le\frac{k+1}{m}\Delta,\,
            \forall j\in[\![1,\ell]\!]
        \right)
        \le
        \exp\left(
            -\frac{d}{\Delta^2}\ell
        \right).
    \end{equation*}
    \bilingual{
        Now for $r_0\le 0\le r_0+\Delta$, we can find some $k\in[\![0,m-1]\!]$ such that $\frac{r_0}\Delta+1\in[\frac km,\frac{k+1}m]$, and hence, for all $\Delta\ge\Delta_0$,
    }{
        现在对 $r_0\le 0\le r_0+\Delta$,~ 存在 $k\in[\![0,m-1]\!]$ 使得 $\frac{r_0}\Delta+1\in[\frac km,\frac{k+1}m]$,~ 因此对所有 $\Delta\ge\Delta_0$,~ 
    }
    \begin{align*}
        &\quad\mathbb P\left(
                r_0\le Y_j\le r_0+\Delta,\,
                \forall j\in[\![1,\ell]\!]
        \right)
        \\ &\le\mathbb P\left(
            \left(-1+\frac km\right)\Delta
            \le Y_j
            \le\frac{k+1}{m}\Delta,\,
            \forall j\in[\![1,\ell]\!]
        \right)
        \le
        \exp\left(
            -\frac{d}{\Delta^2}\ell
        \right).
    \end{align*}
    
    \bilingual{
        Now for any $L\ge 1$, split the interval $[\![1,L]\!]$ into $q$ possible intervals of length $\ell$ and one possible interval of length $r$, i.e., $L=q\ell+r$ with $q\ge 0$ and $0\le r<\ell$. Denote in this proof
    }{
        现在对任何 $L\ge 1$,~ 将区间 $[\![1,L]\!]$ 分成 $q$ 个长度为 $\ell$ 的区间和一个长度为 $r$ 的区间,即 $L=q\ell+r$,~ 其中 $q\ge 0$ 且 $0\le r<\ell$. ~在本证明中记
    }
    \begin{equation*}
        A_k
        :=
        \{
            r_0\le Y_j\le r_0+\Delta,\,
            \forall j\in[\![(k-1)\ell+1,k\ell]\!]
        \},
        \quad
        k\in[\![1,q]\!].
    \end{equation*}
    \bilingual{
        We have
    }{
        从而
    }
    \begin{align*}
        &\quad
        \mathbb P(
            r_0\le Y_j\le r_0+\Delta,\,
            \forall j\in[\![1,L]\!]
        )
        \\ &
        \le
        \mathbb P(
            \cap_{k=1}^qA_k
        )
        =
        \prod_{k=1}^q\mathbb P\left(
            A_k
            \middle|
            \cap_{l=1}^{k-1}A_l
        \right)
        \\ &
        =
        \prod_{k=1}^q\left(
            \frac 1{
                \mathbb P\left(
                    \cap_{l=1}^{k-1}A_l
                \right)
            }
            \mathbb E\left[
                \mathbb E\left(
                    {\bf 1}_{A_k}
                    \middle|
                    \mathscr F_{(k-1)\ell}
                \right)
                {\bf 1}_{
                    \cap_{l=1}^{k-1}A_l
                }
            \right]
        \right)
        \\ &
        \le
        \prod_{k=1}^q
        \sup_{x\in r_0+[0,\Delta]}
        \mathbb P(
            r_0\le Y_j+x\le r_0+\Delta,\,
            \forall j\in[\![1,\ell]\!]
        )
        \\ &
        \le
        \left[
            \exp\left(
                -\frac d{\Delta^2}\ell
            \right)
        \right]^q
        =
        \exp\left(\frac{d}{\Delta^2}r-\frac{d}{\Delta^2}L\right)
        % \\ &
        \le
        \exp\left(\Delta^v-\frac{d}{\Delta^2}L\right),
    \end{align*}
    \bilingual{
        where $r_0-x\in[-\Delta,0]$ for $x\in r_0+[0,\Delta]$, and the last inequality follows from $r<\ell\le\frac{1}{d}\Delta^{2+v}$.
    }{
        其中 $r_0-x\in[-\Delta,0]$ 对 $x\in r_0+[0,\Delta]$ 成立,且最后一式是因 $r<\ell\le\frac{1}{d}\Delta^{2+v}$.
    }
\end{proof}

\begin{remark}
    \bilingual{
        For $u>0$, if we take $\ell=\ell(\Delta)=\lfloor\frac ud\Delta^{2+v}\rfloor$ in lieu of $\ell$ in the previous proof, we can get an upper bound of
    }{
        对 $u>0$,~ 如果在前面的证明中将 $\ell$ 替换为 $\lfloor\frac ud\Delta^{2+v}\rfloor$,~ 可得上界
    }
    $
        e^{
            u\Delta^v-\frac{d}{\Delta^2}L
        }.
    $
\end{remark}

\bilingual{
    Based on the previous lemma, we have the following lemma.
}{
    基于前一个引理,我们有下面的引理.
}

\begin{lemma}[\bilingual{Small Deviation under No Large Drop}{无大回撤意义下的小偏差}, A\"id\'ekon--Hu \mycite{aïdékon2025averageweightpercolationcompletegraph}]
    \label{lemma:2.3}
    \bilingual{
        Let $c>0$. For any $d\in(0,\frac{\pi^2}2\sigma^2)$, there exists $\Delta_0=\Delta_0(c,d)\ge 1$ such that for all $\Delta\ge\Delta_0$ and $L\ge\Delta^{2+c}$, we have
    }{
        设 $c>0$. ~对任何 $d\in(0,\frac{\pi^2}2\sigma^2)$,~ 存在 $\Delta_0=\Delta_0(c,d)\ge 1$ 使得对所有 $\Delta\ge\Delta_0$ 和 $L\ge\Delta^{2+c}$,~ 成立
    }
    {
        \setlength\belowdisplayskip{-20pt}
        \begin{equation*}
            \mathbb P\left(
                Y_L
                \le
                \frac\Delta c,\,
                \min_{0\le i\le j\le L}
                (Y_j-Y_i)
                \ge
                -\Delta
            \right)
            \le
            e^{-\frac{dL}{\Delta^2}}.
        \end{equation*}
    }
\end{lemma}

\begin{figure}[htbp]
    \centering
    \begin{tikzpicture}[
        xscale=1.7, % 修改这里控制横向（代数 n）的间距
        yscale=1.1, % 修改这里控制纵向（空间位置 V）的伸展程度
        >=Stealth,
        % 密集随机游走样式
        rw-path/.style={thick, blue!70!black},
        % 动态屏障样式
        barrier-line/.style={thick, red!70!black, dashed},
        % 参考线样式
        ref-line/.style={thin, gray!60, dashed}
    ]

        % 0. 绘制填充阴影 (放在最底层)
        \fill[blue!5] (0,0) 
            -- (0.2, 0.6) -- (0.4, -0.3) -- (0.6, 0.2) -- (0.8, 1.1) 
            -- (1.1, 0.3) -- (1.3, 0.7) -- (1.6, -0.5) -- (1.9, 0.2) 
            -- (2.2, 1.5) -- (2.5, 0.4) -- (2.8, 0.9) -- (3.1, -0.45) 
            -- (3.4, 0.1) -- (3.7, -0.4) -- (4.0, 0.6) -- (4.3, -0.2) -- (4.5, 0.3) % 蓝色路径结束
            -- (4.5, -0.5) -- (2.2, -0.5) -- (2.1, -0.9) -- (0.8, -0.9) -- (0.7, -1.4) -- (0.2, -1.4) -- (0, -2.0) % 红色路径逆序
            -- cycle;

        % 1. 坐标轴
        \draw[->, thick] (0, -2.5) -- (0, 2.2) node[above] {$Y_n$};
        \draw[->, thick] (0, 0) -- (5.5, 0) node[right] {$n$};

        % 2. 终点限制线 y = Delta/c (假设 Delta/c = 0.8)
        \draw[thick, black!80] (5, 0.8) -- (5.4, 0.8) node[right] {$\frac{\Delta}{c}$};
        \draw[ref-line] (0, 0.8) -- (5, 0.8);

        % 3. 时间 L 的垂线
        \draw[dashed, black!80] (4.5, -2.5) -- (4.5, 2.0) node[above] {$L \ge \Delta^{2+c}$};

        % 4. 密集随机游走路径 (零均值特征)
        % 定义一组波动点
        \draw[rw-path] (0,0) 
            -- (0.2, 0.6) -- (0.4, -0.3) -- (0.6, 0.2) -- (0.8, 1.1) 
            -- (1.1, 0.3) -- (1.3, 0.7) -- (1.6, -0.5) -- (1.9, 0.2) 
            -- (2.2, 1.5) -- (2.5, 0.4) -- (2.8, 0.9) -- (3.1, -0.45) 
            -- (3.4, 0.1) -- (3.7, -0.4) -- (4.0, 0.6) -- (4.3, -0.2) -- (4.5, 0.3);

        % 5. 动态屏障 Max Y_i - Delta (假设 Delta = 2.0)
        % 初始 Max=0 -> 屏障 -2.0
        % 第一高峰 0.6 -> 屏障 -1.4
        % 第二高峰 1.1 -> 屏障 -0.9
        % 第三高峰 1.5 (最高) -> 屏障 -0.5
        \draw[barrier-line] 
            (0, -2.0) -- (0.2, -1.4) -- (0.7, -1.4) -- (0.8, -0.9) 
            -- (2.1, -0.9) -- (2.2, -0.5) -- (4.5, -0.5) 
            node[right] {$\max_{i\le n} Y_i - \Delta$};

        % 6. 标注 Delta 间距 (在最高点处)
        \draw[ref-line] (0, 1.5) -- (2.2, 1.5);
        \draw[<->] (2.3, 1.5) -- (2.3, -0.5) node[midway, left] {$\Delta$};

        % 7. 关键点标注
        \node[below left] at (0,0) {$0$};
        \fill[blue!70!black] (4.5, 0.3) circle (1.5pt) node[right, black] {$Y_L$};

    \end{tikzpicture}
    \caption{\bilingual{Illustration of the historical drawdown of the random walk \texorpdfstring{$Y_n$}{Yn} always constrained by \texorpdfstring{$\Delta$}{Delta} and the endpoint below \texorpdfstring{$\frac{\Delta}{c}$}{Delta/c}}{随机游走 \texorpdfstring{$Y_n$}{Yn} 的历史回撤始终受限于 \texorpdfstring{$\Delta$}{Delta} 且终点位于 \texorpdfstring{$\frac{\Delta}{c}$}{Delta/c} 之下}}
\end{figure}

\begin{proof}
    \label{proof:lemma:2.3}\quad

    \bilingual{
        We denote
    }{
        记事件
    }
    \begin{equation*}
        B = B(c,\Delta,L) := \left\{ Y_L \le \frac\Delta c, \, \min_{0\le i\le j\le L} (Y_j-Y_i) \ge -\Delta \right\}.
    \end{equation*}
    \bilingual{
        Note that $B(c_1,\Delta,L) \subset B(c_2,\Delta,L)$ for all $c_1 \ge c_2$ and fixed $\Delta, L$.
    }{
        注意到对于固定的 $\Delta$ 和 $L$,~ 若 $c_1 \ge c_2$,~ 则有 $B(c_1,\Delta,L) \subset B(c_2,\Delta,L)$.
    }

    \bilingual{
        For the case $c \ge \frac{1}{2}$, it suffices to prove the result for $c' \in (0, \frac{1}{2})$ and $L \ge \Delta^{\frac{2}{1-2c'}}$. Indeed, for any $c \ge \frac{1}{2}$, there exists $c'$ such that
    }{
        对于 $c \ge \frac{1}{2}$ 的情形,只需证明原结论对 $c' \in (0, \frac{1}{2})$ 且 $L \ge \Delta^{\frac{2}{1-2c'}}$ 成立.事实上,对任何 $c \ge \frac{1}{2}$,~ 总能选取 $c'$ 满足
    }
    \begin{equation*}
        0 < \frac{1}{2} - \frac{1}{2+c} \le c' < \frac{1}{2} \le c.
    \end{equation*}
    \bilingual{
        Applying the result for $c'$ and $L \ge \Delta^{\frac{2}{1-2c'}}$, there exists $\Delta_1 = \Delta_1(c',d) = \Delta_1(c,d) \ge 1$ such that for all $\Delta \ge \Delta_1$ and $L \ge \Delta^{\frac{2}{1-2c'}}$
    }{
        对此 $c'$ 和 $L \ge \Delta^{\frac{2}{1-2c'}}$ 应用结论,存在 $\Delta_1 = \Delta_1(c',d) = \Delta_1(c,d) \ge 1$ 使得对所有 $\Delta \ge \Delta_1$ 且 $L \ge \Delta^{\frac{2}{1-2c'}}$ 成立
    }
    \begin{equation*}
        \mathbb{P}(B(c,\Delta,L)) \le \mathbb{P}(B(c',\Delta,L)) \le e^{-\frac{dL}{\Delta^2}}.
    \end{equation*}
    \bilingual{
        Since $\Delta^{\frac{2}{1-2c'}} \le \Delta^{2+c}$, the above inequality holds for all $L \ge \Delta^{2+c}$. In this regime, for $0 < d < \frac{\pi^2}{2}\sigma^2$, we choose $\eta \in (0,1)$ sufficiently small such that
    }{
        由于 $\Delta^{\frac{2}{1-2c'}} \le \Delta^{2+c}$,~ 上式对所有 $L \ge \Delta^{2+c}$ 均成立.此时,对于 $0 < d < \frac{\pi^2}{2}\sigma^2$,~ 选取足够小的 $\eta \in (0,1)$ 使得
    }
    \begin{equation*}
        (2+\eta)\left(\frac{1}{2}-c\right) < 1, \quad d(1+2\eta)^4 < \frac{\pi^2}{2}\sigma^2.
    \end{equation*}

    \bilingual{
        For $0 < c < \frac{1}{2}$, we directly prove the result for $L \ge \Delta^{2+c}$ using the same method, only requiring $\eta \in (0,1)$ to be small enough to satisfy $\eta < c$ and the second inequality above.
    }{
        对于 $0 < c < \frac{1}{2}$ 的情形,直接用相同方法证明 $L \ge \Delta^{2+c}$ 时的结论,区别仅在于需令 $\eta \in (0,1)$ 足够小以满足 $\eta < c$ 且成立上述第二个不等式.
    }

    \bilingual{
        Note that in both cases, we have
    }{
        注意到,无论上述哪种情况,均满足
    }
    \begin{equation*}
        \frac{1}{c} \ge 2, \quad L > \lfloor\Delta^{2+\eta}\rfloor, \quad \lim_{\Delta\to\infty} \frac{L}{\lfloor\Delta^{2+\eta}\rfloor} = \infty.
    \end{equation*}
    \bilingual{
        We divide the interval $[\![0,L]\!]$ into smaller intervals of length $\lfloor\Delta^{2+\eta}\rfloor$, and the space into intervals of length $\lceil\eta\Delta\rceil$. Let $q := \lfloor\frac{L}{\lfloor\Delta^{2+\eta}\rfloor}\rfloor$, $\ell_0 := 0$, and
    }{
        将时间区间 $[\![0,L]\!]$ 分成长度为 $\lfloor\Delta^{2+\eta}\rfloor$ 的小区间,并将空间划分为长度为 $\lceil\eta\Delta\rceil$ 的层级.令 $q := \lfloor\frac{L}{\lfloor\Delta^{2+\eta}\rfloor}\rfloor$,~ $\ell_0 := 0$,~ 且定义
    }
    \begin{align*}
        \ell_j &:= L - q\lfloor\Delta^{2+\eta}\rfloor + (j-1)\lfloor\Delta^{2+\eta}\rfloor, \quad j \in [\![1, q+1]\!], \\
        I_j &:= [\![\ell_j, \ell_{j+1}]\!], \quad j \in [\![0, q]\!].
    \end{align*}
    \bilingual{
        On the event $B$, we have
    }{
        在事件 $B$ 上,成立
    }
    \begin{equation*}
        \min_{0\le j\le L} Y_j \ge -\Delta, \quad \max_{0\le i\le L} Y_i \le Y_L + \Delta \le \left(\frac{1}{c}+1\right)\Delta.
    \end{equation*}
    \bilingual{
        Let $C := \frac{1}{c}+2 \ge 4$ so that the total range is within $C\Delta$. For any $j \in [\![0,q]\!]$, there exists some random $k_j \in \mathbb{N}$ such that
    }{
        令 $C := \frac{1}{c}+2 \ge 4$,~ 此时路径取值的跨度不超过 $C\Delta$. ~对于任何 $j \in [\![0,q]\!]$,~ 存在随机变量 $k_j \in \mathbb{N}$ 满足
    }
    \begin{equation*}
        \left(\frac{1}{c}+1\right)\Delta - (k_j+1)\lceil\eta\Delta\rceil < \max_{0 \le \ell \le \ell_{j+1}} Y_\ell \le \left(\frac{1}{c}+1\right)\Delta - k_j\lceil\eta\Delta\rceil := M_j.
    \end{equation*}
    \bilingual{
        Note that $M_j$ can take at most $\lceil \frac C\eta \rceil$ possible values. Let $\mathcal{M}$ be the deterministic set of these possible values. Define the index set
    }{
        注意 $M_j$ 的可能取值个数不超过 $\lceil \frac C\eta \rceil$ 个.令 $\mathcal{M}$ 为 $M_j$ 可能取值的确定性集合.定义下标集
    }
    \begin{equation*}
        J := \left\{ j \in [\![0,q-1]\!] : M_j \neq M_{j+1} \right\}.
    \end{equation*}
    \bilingual{
        Since $M_j$ is non-decreasing in $j$, we have $\#J \le \lceil \frac C\eta \rceil$.
    }{
        由于 $M_j$ 随 $j$ 非递减,故有 $\#J \le \lceil \frac C\eta \rceil$.
    }

    \bilingual{
        For $j \notin J$, we have $M_j = M_{j+1}$, which implies that on $B$, $Y_i \in [M_j - \lceil\eta\Delta\rceil - \Delta, M_j]$ for all $i \in I_{j+1}$. (See \Cref{figure:tube-constraint}.)
    }{
        若 $j \notin J$,~ 则 $M_j = M_{j+1}$,~ 由此推得,在事件 $B$ 上,对所有 $i \in I_{j+1}$,~ 都有 $Y_i \in [M_j - \lceil\eta\Delta\rceil - \Delta, M_j]$. (见\Cref{figure:tube-constraint}.)
    }

    \begin{figure}[H]
        \centering
        \begin{tikzpicture}[
            xscale=1.2, 
            yscale=1.2, 
            >=Stealth,
            rw-path/.style={thick, blue!70!black},
            running-max/.style={thin, gray!80, dashed},
            barrier-red/.style={thick, red!70!black, dashed},
            barrier-green/.style={thick, green!60!black, dashed},
            ref-line/.style={thin, black!80, dashed}
        ]

            % 1. 参数设定
            % Delta = 2.2
            % Mj = 1.8, Mj - etaDelta = 1.2
            % 绿线 = 1.2 - 2.2 = -1.0
            
            % 2. 绘制阴影区域 (Y_n 与 红色动态屏障之间)
            \fill[blue!5] 
                (0,0) -- (0.5, 1.1) -- (1.5, -0.8) -- (2.5, 1.6) -- (3.5, 0.4) -- (4, 1.5) % I_j 结束
                -- (5, -0.5) -- (6, 1.7) -- (7, -0.2) -- (8, 0.5) % I_{j+1} 结束
                -- (8, -0.5) -- (6, -0.5) -- (5.95, -0.6)-- (2.5, -0.6) -- (2.3, -1.1) -- (0.5, -1.1) -- (0, -2.2) % 红色路径逆序
                -- cycle;

            % 3. 坐标轴
            % \draw[->, thick] (-0.2, -2.5) -- (8.5, -2.5) node[right] {$n$};
            \draw[->, thick] (-0.2, -2.5) -- (-0.2, 2.5) node[above] {$Y_n$};

            % 4. 离散层级标注 (仅标注这两根线)
            \draw[barrier-green] (-0.3, 1.8) -- (8.2, 1.8) node[right] {$M_j$};
            \draw[black!80, dashed] (-0.3, 1.2) -- (8.2, 1.2) node[right] {$M_j - \lceil\eta\Delta\rceil$};

            % 5. 时间区间划分
            \draw[ref-line] (0, -2.6) -- (0, 2.2) node[below=5.6cm] {$\ell_j$};
            \draw[ref-line] (4, -2.6) -- (4, 2.2) node[below=5.6cm] {$\ell_{j+1}$};
            \draw[ref-line] (8, -2.6) -- (8, 2.2) node[below=5.6cm] {$\ell_{j+2}$};
            \node at (2, -2.8) {$I_j$};
            \node at (6, -2.8) {$I_{j+1}$};

            % 6. 静态绿线 (区间起始限制)
            \draw[barrier-green] (4, -1.0) -- (8.2, -1.0) node[right] {$(M_j - \lceil\eta\Delta\rceil) - \Delta$};

            % 7. 零均值随机游走路径
            \draw[rw-path] (0,0) 
                -- (0.5, 1.1) -- (1.5, -0.8) -- (2.5, 1.6) -- (3.5, 0.4) -- (4, 1.5) 
                -- (5, -0.5) -- (6, 1.7) -- (7, -0.2) -- (8, 0.5);

            % 8. 区间最大值(Running Max 阶梯线)
            \draw[running-max] (0,0) -- (0.5, 1.1) -- (2.3, 1.1) -- (2.5, 1.6) -- (5.95, 1.6) -- (6, 1.7) -- (8, 1.7);

            % 9. 动态红线 (max Y_i - Delta)
            % 随创新高而阶梯式上升
            \draw[barrier-red] (0, -2.2) -- (0.5, -1.1) -- (2.3, -1.1) -- (2.5, -0.6) -- (5.95, -0.6) -- (6, -0.5) -- (8, -0.5)
                node[right=0.25cm] {$\max_{i\le n} Y_i - \Delta$};

            % 10. 标注 Delta 间距 (在 I_{j+1} 内部展示)
            \draw[<->] (6.2, 1.7) -- (6.2, -0.5);
            \node[right] at (6.2, 0.6) {$\Delta$};

            % 11. 标注起点 0
            % \node[left] at (0,0) {$0$};

        \end{tikzpicture}
        \caption{\bilingual{Illustration of the tube constraints for the random walk when \texorpdfstring{$j\notin J$}{j not in J} within intervals \texorpdfstring{$I_j$}{Ij} and \texorpdfstring{$I_{j+1}$}{Ij+1}}{\texorpdfstring{$j\notin J$}{j not in J}时区间 \texorpdfstring{$I_j$}{Ij} 与 \texorpdfstring{$I_{j+1}$}{Ij+1} 内随机游走的管道约束}}
        \label{figure:tube-constraint}
    \end{figure}

    \bilingual{
        Now, let $\mathscr{J}$ be the family of all possible indices where the discretized maximum can jump
    }{
        令 $\mathscr{J}$ 为各区间最大值所有可能跳跃点的下标集族
    }
    \begin{equation*}
        \mathscr{J} := \left\{ J_0 \subseteq [\![0, q-1]\!] : \#J_0 \le \left\lceil \frac{C}{\eta} \right\rceil \right\}.
    \end{equation*}
    \bilingual{
        For a fixed $J_0 \in \mathscr{J}$, let its complement be $J_0^{\mathrm{c}} = \{ j_1 > j_2 > \cdots > j_{\#J_0^{\mathrm{c}}} \} := [\![0, q-1]\!] \setminus J_0$. Note that $\#J_0^{\mathrm{c}} \ge q - \lceil \frac C\eta \rceil$.
    }{
        对于固定的 $J_0 \in \mathscr{J}$,~ 令其补集为 $J_0^{\mathrm{c}} = \{ j_1 > j_2 > \cdots > j_{\#J_0^{\mathrm{c}}} \} := [\![0, q-1]\!] \setminus J_0$. ~注意 $\#J_0^{\mathrm{c}} \ge q - \lceil \frac C\eta \rceil$.
    }
    
    \bilingual{
        We estimate $\mathbb{P}(B, J=J_0)$ by applying the Markov property at times $\ell_{j_k+1}$ in decreasing order of $k$. Since $B \cap \{J=J_0\}$ implies that for all $j \in J_0^{\mathrm{c}}$, the maximum $M_j$ remains constant over the interval $I_{j+1}$, the path is constrained within a tube of width $\lceil\eta\Delta\rceil + \Delta$. We obtain
    }{
        通过按 $k$ 降序对时刻 $\ell_{j_k+1}$ 应用马尔可夫性质来估计 $\mathbb{P}(B, J=J_0)$. ~由于事件 $B \cap \{J=J_0\}$ 蕴含了:对所有 $j \in J_0^{\mathrm{c}}$,~ 区间最大值 $M_j$ 在区间 $I_{j+1}$ 上保持不变,此时路径被限制在宽度为 $\lceil\eta\Delta\rceil + \Delta$ 的管道内.由此得
    }
    \begin{align*}
        \mathbb{P}(B, J=J_0) 
        &\le \mathbb{E} \left[ \prod_{k=1}^{\#J_0^{\mathrm{c}}} \mathbf{1}_{ \left\{ \forall i \in I_{j_k+1}, \, M_{j_k}-\lceil\eta\Delta\rceil-\Delta \le Y_i \le M_{j_k} \right\} } \right] \\
        &= \mathbb{E} \left[ \prod_{k=2}^{\#J_0^{\mathrm{c}}} \mathbf{1}_{\{\dots\}} \mathbb{E} \left( \mathbf{1}_{ \left\{ \forall i \in I_{j_1+1}, \, M_{j_1}-\lceil\eta\Delta\rceil-\Delta \le Y_i \le M_{j_1} \right\} } \middle| \mathcal{F}_{\ell_{j_1+1}} \right) \right].
    \end{align*}
    \bilingual{
        Since $M_{j_1}$ is $\mathcal{F}_{\ell_{j_1+1}}$-measurable, we can evaluate the inner conditional expectation by considering the starting position $a = Y_{\ell_{j_1+1}}$. On the event $B$, $a$ is constrained within $[M_{j_1}-\lceil\eta\Delta\rceil-\Delta, M_{j_1}]$. Thus, for any $m \in \mathcal{M}$, we have
    }{
        由于 $M_{j_1}$ 是 $\mathcal{F}_{\ell_{j_1+1}}$-可测的,我们可以通过考察起始位置 $a = Y_{\ell_{j_1+1}}$ 来计算内部条件期望.在事件 $B$ 上,~ $a$ 被限制在 $[M_{j_1}-\lceil\eta\Delta\rceil-\Delta, M_{j_1}]$ 内.因此,对于任何 $m \in \mathcal{M}$,~ 成立
    }
    \begin{align*}
        &\quad \mathbb{E} \left( \mathbf{1}_{\{\dots\}} \middle| \mathcal{F}_{\ell_{j_1+1}} \right) \\
        &=\sum_{m\in\mathcal M} \mathbf{1}_{\left\{M_{j_1}=m\right\}} \mathbb{P} \left.\left( \forall i \in [\![1, \lfloor\Delta^{2+\eta}\rfloor]\!], \, m-\lceil\eta\Delta\rceil-\Delta \le Y_i + a \le m \right)\right|_{a=Y_{\ell_{j_1+1}}} \\
        &\le\sum_{m\in\mathcal M} \mathbf{1}_{\left\{M_{j_1}=m\right\}} \sup_{a \in [m-\lceil\eta\Delta\rceil-\Delta, m]} \mathbb{P} \left( \forall i \in [\![1, \lfloor\Delta^{2+\eta}\rfloor]\!], \, m-\lceil\eta\Delta\rceil-\Delta \le Y_i + a \le m \right) \\
        &\le \max_{m \in \mathcal{M}} \sup_{a \in [m-\lceil\eta\Delta\rceil-\Delta, m]} \mathbb{P} \left( \forall i \in [\![1, \lfloor\Delta^{2+\eta}\rfloor]\!], \, m-\lceil\eta\Delta\rceil-\Delta \le Y_i + a \le m \right) \\
        &\le e^{
            (\lceil\eta\Delta\rceil+\Delta)^{\eta/2}
            -
            \frac{
                d(1+2\eta)^4
            }{
                (\lceil\eta\Delta\rceil+\Delta)^2
            }
            \lfloor\Delta^{2+\eta}\rfloor
        }
        \le
        e^{
            (2\eta\Delta+\Delta)^{\eta/2}
            -
            \frac{
                d(1+2\eta)^4
            }{
                (2\eta\Delta+\Delta)^2
            }
            (\Delta^{2+\eta}-\Delta^{2+\eta/2})
        }
        \\ &
        \le
        e^{
            2d\eta(1+2\eta)\Delta^\eta
            -
            d(1+2\eta)^2\Delta^\eta
        }
        =
        e^{
            -d(1+2\eta)\Delta^\eta
        },
    \end{align*}
    \bilingual{
        where the fourth inequality follows from \Cref{lemma:4.1} applied to a tube of width $\lceil\eta\Delta\rceil+\Delta$. This holds for some $\Delta_1=\Delta_1(c,d)>0$ and all $\lceil\eta\Delta\rceil+\Delta\ge\Delta_1$, with parameters $v=\frac\eta 2$, $d(1+2\eta)^4$ in place of $d$, and $\lfloor\Delta^{2+\eta}\rfloor$ in place of $L$. The final inequality holds for some $\Delta_2=\Delta_2(c,d)>\frac 1{1+\eta}\Delta_1$ and all $\Delta\ge\Delta_2$. Iterating this argument for all $k$ yields
    }{
        其中第四个不等式对宽度为 $\lceil\eta\Delta\rceil+\Delta$ 的管道应用\Cref{lemma:4.1} 得到,对于某个 $\Delta_1=\Delta_1(c,d)>0$ 和所有 $\lceil\eta\Delta\rceil+\Delta\ge\Delta_1$ 均成立.参数设置为 $v=\frac\eta 2$,~ 用 $d(1+2\eta)^4$ 替代 $d$,~ 且用 $\lfloor\Delta^{2+\eta}\rfloor$ 替代 $L$. ~而最后一个不等式对某个 $\Delta_2=\Delta_2(c,d)>\frac 1{1+\eta}\Delta_1$ 和所有 $\Delta\ge\Delta_2$ 成立.对所有 $k$ 迭代此过程可得
    }
    \begin{align*}
        \mathbb P(B,J=J_0)
        &
        \le
        \left( e^{ -d(1+2\eta)\Delta^\eta } \right)^{\#J_0^{\mathrm c}}
        \le
        e^{
            -d(1+2\eta)\Delta^\eta
            \left(
                \lfloor \frac L{\lfloor\Delta^{2+\eta}\rfloor} \rfloor
                -
                \lceil \frac C\eta \rceil
            \right)
        }
        \\ &
        \le
        e^{
            -d(1+2\eta)\Delta^\eta
            \left(
                \frac L{\Delta^{2+\eta}}
                -1
                -
                2\frac C\eta
            \right)
        }
        =
        e^{
            -d
            (1+2\eta)
            \frac L{\Delta^2}
            + d
            (1+\eta)(1+2\frac C\eta)
            \Delta^\eta
        }
        \\ &
        \le
        e^{
            -d
            (1+2\eta)
            \frac L{\Delta^2}
            + d
            \eta
            \frac L{\Delta^2}
        }
        =
        e^{
            -d
            \left(1+\eta\right)
            \frac L{\Delta^2}
        },
    \end{align*}
    \bilingual{
        where the last inequality holds for some $\Delta_3=\Delta_3(c,d)>\Delta_2$ and all $\Delta\ge\Delta_3$. Finally, summing over all $J_0 \in \mathscr{J}$, we have
    }{
        其中最后一个不等式对某个 $\Delta_3=\Delta_3(c,d)>\Delta_2$ 和所有 $\Delta\ge\Delta_3$ 成立.于是
    }
    \begin{align*}
        \mathbb P(B)
        &
        =
        \sum_{J_0\in\mathscr J}
        \mathbb P(B,J=J_0)
        \le
        \binom{q}{\lceil C/\eta \rceil}
        e^{
            -d
            \left(1+\eta\right)
            \frac L{\Delta^2}
        }
        \le
        2^{q}
        e^{
            -d
            \left(1+\eta\right)
            \frac L{\Delta^2}
        }
        \\ &
        \le
        e^{
            \frac L{\Delta^{2+\eta/2}}\log 2
            -d
            (1+\eta)
            \frac L{\Delta^2}
        }
        \le
        e^{
            d
            \eta
            \frac L{\Delta^2}
            -d
            (1+\eta)
            \frac L{\Delta^2}
        }
        =
        e^{
            -d
            \frac L{\Delta^2}
        },
    \end{align*}
    \bilingual{
        where the inequalities hold for some $\Delta_0=\Delta_0(c,d)>\Delta_3$ and all $\Delta\ge\Delta_0$.
    }{
        其中不等式对某个 $\Delta_0=\Delta_0(c,d)>\Delta_3$ 和所有 $\Delta\ge\Delta_0$ 成立,证毕.
    }
\end{proof}

\bilingual{
    Based on \Cref{lemma:2.3}, \Cref{lemma:4.1} can be extended to include the case of a small drift.
}{
    基于\Cref{lemma:2.3},~ \Cref{lemma:4.1} 可以推广到包含小漂移的情形.
}

\begin{lemma}[\bilingual{Small Deviation with Small Drift}{小漂移倾斜管道意义下的小偏差}]\label{lemma:4.2}
    \bilingual{
        Let $v>0$, $0<d<\frac{\pi^2}2\sigma^2$ and $c\ge 0$. There exists $\Delta_0>0$ such that for all $\Delta\ge\Delta_0$, $r_0\le 0\le r_0+\Delta$ and $L\ge 1$,
    }{
        设 $v>0$,~ $0<d<\frac{\pi^2}2\sigma^2$ 且 $c\ge 0$. ~存在 $\Delta_0>0$ 使得对所有 $\Delta\ge\Delta_0$,~ $r_0\le 0\le r_0+\Delta$ 且 $L\ge 1$,
    }
    \begin{equation}
        \mathbb P\left(
            r_j\le Y_j\le s_j,\,
            \forall j\in[\![1,L]\!]
        \right)
        \le
        \exp\left(
            \Delta^v-\frac d{\Delta^2}L
        \right),
        \label{eq:4.2}
    \end{equation}
    \bilingual{
        where $r_j:=r_0+\frac c{\Delta^2}j$ and $s_j:=r_j+\Delta$.
    }{
        其中 $r_j:=r_0+\frac c{\Delta^2}j$ 且 $s_j:=r_j+\Delta$.
    }
\end{lemma}

\begin{figure}[htbp]
    \centering
    \begin{tikzpicture}[
        xscale=1.5,
        yscale=1.0,
        >=Stealth,
        rw-path/.style={thick, blue!70!black, mark=*, mark size=1.5pt},
        tube-fill/.style={gray!10},
        label-font/.style={font=\small}
    ]

        % 1. 计算坐标参数 (斜率为 0.2)
        % x=0: r0 = -1.5, s0 = 0.5
        % x=4: r4 = -1.5 + 0.2*4 = -0.7, s4 = 0.5 + 0.2*4 = 1.3
        % x=4.5: r4.5 = -1.5 + 0.2*4.5 = -0.6, s4.5 = 0.5 + 0.2*4.5 = 1.4

        % 2. 背景倾斜填充 (平行四边形区域)
        \fill[tube-fill] (0, -1.5) -- (4, -0.7) -- (4, 1.3) -- (0, 0.5) -- cycle;

        % 3. 参考虚线
        \draw[dashed, black!60] (0, -2) -- (4.5, -2) node[right] {$-\Delta$};
        \draw[dashed, black!60] (4, -2.5) -- (4, 1.8) node[above] {$L$};

        % 4. 倾斜管道边界 (实线)
        % 下边界: y = 0.2x - 1.5
        \draw[thick, black] (0, -1.5) -- (4.5, -0.6) node[right] {$r_n$};
        % 上边界: y = 0.2x + 0.5
        \draw[thick, black] (0, 0.5) -- (4.5, 1.4) node[right] {$s_n$};

        % 5. 坐标轴 (n轴虚线化，突出零基准)
        \draw[->, thick] (0, -2.5) -- (0, 2) node[above] {$Y_n$};
        \draw[->, thick, dashed] (0, 0) -- (5, 0) node[right] {$n$};

        % 6. 随机游走路径 (根据斜率向上微调)
        \draw[rw-path] plot coordinates {
            (0, 0) 
            (0.8, -0.3) 
            (1.6, 0.6) 
            (2.4, -0.1) 
            (3.2, 0.8) 
            (4, 0.4)
        };

        % 7. 左侧跨度标注 (x=0)
        \draw[decorate, decoration={brace, amplitude=5pt}] 
            (-0.1, -1.5) -- (-0.1, 0.5) node[midway, left=6pt] {$\Delta$};

        % 8. 右侧跨度标注 (x=l 的右边)
        \draw[decorate, decoration={brace, amplitude=5pt, mirror}] 
            (4.1, -0.7) -- (4.1, 1.3) node[midway, right=6pt] {$\Delta$};

        % 标注起始点 r0
        \node[left] at (0, -1.5) {$r_0$};

    \end{tikzpicture}
    \caption{\bilingual{Illustration of the small deviation of the random walk \texorpdfstring{$Y_n$}{Yn} within the tilted tube \texorpdfstring{$[r_n, s_n]$}{[r_n, s_n]} with small drift}{随机游走 \texorpdfstring{$Y_n$}{Yn} 在具有小漂移的倾斜管道 \texorpdfstring{$[r_n, s_n]$}{[r_n, s_n]} 内的小偏差示意图}}
\end{figure}

\begin{proof}\label{proof:lemma:4.2}
    \quad

    \bilingual{
        Fix $0<\eta<v\wedge 1$. Let $\ell=\ell(\Delta):=\lfloor\Delta^{2+\eta}\rfloor$. We first note that if $L<\ell$, then for some $\Delta_1=\Delta_1(c,d,v)>0$ and all $\Delta\ge\Delta_1$, we have $d\Delta^\eta<\Delta^v$ since $\eta<v$. Hence, the right-hand side of \eqref{eq:4.2} is greater than $1$ and the inequality trivially holds.
    }{
        固定 $0<\eta<v\wedge 1$. ~令 $\ell=\ell(\Delta):=\lfloor\Delta^{2+\eta}\rfloor$. ~首先注意到如果 $L<\ell$,~ 则对于某个 $\Delta_1=\Delta_1(c,d,v)>0$ 和所有 $\Delta\ge\Delta_1$,~ 成立 $d\Delta^\eta<\Delta^v$, 此因 $\eta<v$. ~因此,~ \eqref{eq:4.2} 的右侧大于 $1$ 并且不等式显然成立.
    }

    \bilingual{
        Now assume $L \ge \ell$. Divide the interval $[\![1, L]\!]$ into $q = \lfloor\frac L\ell \rfloor \ge 1$ intervals of length $\ell$, and a remainder interval of length $r = L - q\ell \in [0, \ell)$. 
        Define the event $B := \{ r_j \le Y_j \le s_j, \, \forall j \in [\![1, L]\!] \}$. 
        For each $k \in [\![1, q]\!]$, let $A_k$ be the event that the random walk stays in the tilted tube during the $k$-th interval, i.e.,
    }{
        现假设 $L \ge \ell$. ~将区间 $[\![1, L]\!]$ 分成 $q = \lfloor \frac L\ell \rfloor \ge 1$ 个长度为 $\ell$ 的区间,以及一个余项区间,长度为 $r = L - q\ell \in [0, \ell)$.
        定义事件 $B := \{ r_j \le Y_j \le s_j, \, \forall j \in [\![1, L]\!] \}$. 
        对于每个 $k \in [\![1, q]\!]$,~ 令 $A_k$ 为随机游走在第 $k$ 个区间内保持在倾斜管道中的事件,即
    }
    $
        A_k := \left\{ r_j \le Y_j \le s_j, \, \forall j \in [\![(k-1)\ell+1, k\ell]\!] \right\}.
    $
    \bilingual{
        Clearly $B \subseteq \bigcap_{k=1}^q A_k$. By the Markov property,
    }{
        显然 $B \subseteq \bigcap_{k=1}^q A_k$. ~利用Markov 性质有
    }
    \begin{equation*}
        \mathbb{P}(B) \le \mathbb{E} \left[ \mathbb{E} \left( \mathbf{1}_{A_q} \middle| \mathcal{F}_{(q-1)\ell} \right) \mathbf{1}_{\bigcap_{k=1}^{q-1} A_k} \right].
    \end{equation*}
    \bilingual{
        For any state $y$ such that $Y_{(q-1)\ell} = y$ on  $\bigcap_{k=1}^{q-1} A_k$, we have $y \in [r_{(q-1)\ell}, s_{(q-1)\ell}]$. The conditional probability is
    }{
        对于事件 $\bigcap_{k=1}^{q-1} A_k$ 上的任何状态 $y = Y_{(q-1)\ell}$,~ 满足 $y \in [r_{(q-1)\ell}, s_{(q-1)\ell}]$. ~其条件概率为
    }
    \begin{align*}
        \mathbb{P}_y(A_q) 
        &= \mathbb{P} \left( r_{(q-1)\ell+j} \le y + Y_j \le s_{(q-1)\ell+j}, \, \forall j \in [\![1, \ell]\!] \right) \\
        &= \mathbb{P} \left( r_{(q-1)\ell} + \textstyle\frac c{\Delta^2}j \le y + Y_j \le r_{(q-1)\ell} + \frac c{\Delta^2}j + \Delta, \, \forall j \in [\![1, \ell]\!] \right) \\
        &= \mathbb{P} \left( r_0' + \textstyle\frac c{\Delta^2}j \le Y_j \le r_0' + \frac c{\Delta^2}j + \Delta, \, \forall j \in [\![1, \ell]\!] \right),
    \end{align*}
    \bilingual{
        where $r_0' := r_{(q-1)\ell} - y$. Since $y \in [r_{(q-1)\ell}, r_{(q-1)\ell} + \Delta]$, we have $r_0' \in [-\Delta, 0]$. 
        On this event, for any $0 \le i < j \le \ell$, the drawdown satisfies $Y_j - Y_i \ge \frac c{\Delta^2}(j-i) - \Delta \ge -\Delta$. 
        Also, the terminal value satisfies $Y_\ell \le r_0' + \frac c{\Delta^2}\ell + \Delta \le \frac c{\Delta^2}\ell + \Delta$. 
        Since $\frac c{\Delta^2}\ell \le \frac{c}{\Delta^2} \Delta^{2+\eta} = c \Delta^\eta$ and $\eta < 1$, we have $Y_\ell \le \Delta (1 + c \Delta^{\eta-1}) \le \frac {2\Delta} \eta$ for some $\Delta_2=\Delta_2(c,d,v)>\Delta_1$ and all $\Delta \ge \Delta_2$. 
        Applying \Cref{lemma:2.3} with $\frac \eta 2$ in lieu of $c$, we have, for some $\Delta_0=\Delta_0(c,d,v)>\Delta_2$ and all $\Delta \ge \Delta_0$,
    }{
        其中 $r_0' := r_{(q-1)\ell} - y$. ~由于 $y \in [r_{(q-1)\ell}, r_{(q-1)\ell} + \Delta]$,~ 成立 $r_0' \in [-\Delta, 0]$. 
        在此事件上,对于任何 $0 \le i < j \le \ell$,~ 回撤满足 $Y_j - Y_i \ge \frac c{\Delta^2}(j-i) - \Delta \ge -\Delta$.
        同时,终点值满足 $Y_\ell \le r_0' + \frac c{\Delta^2}\ell + \Delta \le \frac c{\Delta^2}\ell + \Delta$.
        由于 $\frac c{\Delta^2}\ell \le \frac{c}{\Delta^2} \Delta^{2+\eta} = c \Delta^\eta$ 且 $\eta < 1$,~ 对于某个 $\Delta_2=\Delta_2(c,d,v)>\Delta_1$ 和所有 $\Delta \ge \Delta_2$,~ 成立 $Y_\ell \le \Delta (1 + c \Delta^{\eta-1}) \le \frac {2\Delta} \eta$. 
        应用\Cref{lemma:2.3},参数设为$\frac\eta 2$代替$c$,~ 可得对于某个$\Delta_0=\Delta_0(c,d,v)>\Delta_2$和所有$\Delta\ge\Delta_0$,~ 有
    }
    \begin{equation*}
        \mathbb{E} \left( \mathbf{1}_{A_q} \middle| \mathcal{F}_{(q-1)\ell} \right) \le \sup_{r_0' \in [-\Delta, 0]} \mathbb{P}_{y} (A_q) \le \exp\left( -d \frac{\ell}{\Delta^2} \right).
    \end{equation*}
    \bilingual{
        By induction, we obtain
    }{
        通过归纳迭代,得到
    }
    \begin{equation*}
        \mathbb{P}(B) \le \prod_{k=1}^q \exp\left( -d \frac{\ell}{\Delta^2} \right) = \exp\left( - \frac{d (q\ell)}{\Delta^2} \right).
    \end{equation*}
    \bilingual{
        Substituting $q\ell = L - r$ and noting that $r<\ell$ and $\Delta\ge\Delta_1$, we have
    }{
        代入 $q\ell = L - r$,~ 注意到$r<\ell$与$\Delta\ge\Delta_1$,~ 有
    }
    \begin{equation*}
        \mathbb{P}(B) \le \exp\left( - \frac{dL}{\Delta^2} + \frac{dr}{\Delta^2} \right) \le \exp\left( - \frac{dL}{\Delta^2} + d\Delta^\eta \right) \le \exp\left( \Delta^v - \frac{dL}{\Delta^2} \right),
    \end{equation*}
    \bilingual{
        which completes the proof.
    }{
        证毕.
    }
\end{proof}

\begin{remark}
    \bilingual{
        For $u>0$, if we take $d\Delta^\eta<u\Delta^v$ in lieu of $d\Delta^\eta<\Delta^v$ in the previous proof, we can get an upper bound of
    }{
        对 $u>0$,~ 在前面的证明中用 $d\Delta^\eta<u\Delta^v$ 代替 $d\Delta^\eta<\Delta^v$,~ 可得上界
    }
    $
        e^{
            u\Delta^v-\frac{d}{\Delta^2}L
        }.
    $
\end{remark}

\bilingual{
    The following lemma, using \nameref{theorem:mogulskii}, is similar to \Cref{lemma:4.1} and provides a lower bound for the small deviation probability under the tilted tube end constraint.
}{
    以下引理类似\Cref{lemma:4.1},利用\nameref{theorem:mogulskii},提供了倾斜管道末端限制下的小偏差下界.
}

\begin{lemma}[\bilingual{Lower Bound for Small Deviation under Tilted Tube End Constraint}{倾斜管道末端限制下的小偏差下界}]\label{lemma:4.3}
    \bilingual{
        Let $v>0$, $d'>\frac{\pi^2}{2}\sigma^2$, $c\ge 0$ and $\varepsilon>0$. There exists $\Delta_0>0$ such that for all $\Delta\ge\Delta_0$, $r_0\le 0\le r_0+\Delta$ and $L\ge\lfloor\Delta^{2+\epsilon}\rfloor$,
    }{
        设 $v>0$,~ $d'>\frac{\pi^2}{2}\sigma^2$,~ $c\ge 0$ 且 $\varepsilon>0$. ~存在 $\Delta_0>0$ 使得对所有 $\Delta\ge\Delta_0$,~ $r_0\le 0\le r_0+\Delta$ 且 $L\ge\lfloor\Delta^{2+\epsilon}\rfloor$,
    }
    \begin{align}
        \mathbb P\left(
            r_j\le Y_j\le s_j,\,
            \forall j\in[\![1,L]\!],\,
            Y_L\in[r_L,r_L+v\Delta]
        \right)
        &\ge
        \exp\left(-\frac{d'}{\Delta^2}L\right),
        \label{eq:4.3}
        \\
        \mathbb P\left(
            r_j\le Y_j\le s_j,\,
            \forall j\in[\![1,L]\!],\,
            Y_L\in[s_L-v\Delta,s_L]
        \right)
        &\ge
        \exp\left(-\frac{d'}{\Delta^2}L\right),
        \label{eq:4.3:2}
    \end{align}
    \bilingual{
        where $r_j=r_0+\frac{c}{\Delta^2}j$ and $s_j=r_j+\Delta$.
    }{
        其中 $r_j=r_0+\frac{c}{\Delta^2}j$ 且 $s_j=r_j+\Delta$.
    }
\end{lemma}

\begin{figure}[htbp]
    \centering
    \begin{tikzpicture}[
        xscale=1.7, % 修改这里控制横向（代数 n）的间距
        yscale=1.1, % 修改这里控制纵向（空间位置 V）的伸展程度
        >=Stealth,
        rw-path/.style={thick, blue!70!black, mark=*, mark size=1.5pt},
        tube-fill/.style={gray!10},
        label-font/.style={font=\small}
    ]

        % 1. 坐标参数定义 (斜率 0.2, v=0.4)
        % x=0: r0 = -1.5, s0 = 0.5
        % x=4: rl = -0.7, sl = 1.3, rl+vDelta = -0.7 + 0.8 = 0.1

        % 2. 背景倾斜填充
        \fill[tube-fill] (0, -1.5) -- (4, -0.7) -- (4, 1.3) -- (0, 0.5) -- cycle;

        % 3. 参考虚线与时间尺度标注
        \draw[dashed, black!60] (0, -2) -- (4.5, -2) node[right] {$-\Delta$};
        \draw[dashed, black!60] (4, -2.5) -- (4, 1.8) node[above, black] {$L \ge \Delta^{2+\varepsilon}$};

        % 4. 倾斜管道边界 (实线)
        \draw[thick, black] (0, -1.5) -- (4.5, -0.6) node[right] {$r_n$};
        \draw[thick, black] (0, 0.5) -- (4.5, 1.4) node[right] {$s_n$};

        % 5. 末端限制平台 (x=l处的小平台)
        \draw[dashed, black] (4, -0.7) -- (6, -0.7);
        \draw[dashed, black] (4, 0.1) -- (6, 0.1);
        % 标注 v*Delta 间距
        \draw[<->, black] (6, -0.7) -- (6, 0.1) node[midway, right] {$v\Delta$};
        \node at (6, -0.7) [below] {$r_L$};

        % 6. 坐标轴
        \draw[->, thick] (0, -2.5) -- (0, 2) node[above] {$Y_n$};
        \draw[->, thick, dashed] (0, 0) -- (5, 0) node[right] {$n$};

        % 7. 密集的随机游走路径 (根据斜率向上微调并落在目标区间 [-0.7, 0.1] 内)
        \draw[rw-path] (0, 0) 
            -- (0.2, 0.3) -- (0.4, -0.2) -- (0.6, 0.1) -- (0.8, -0.4) 
            -- (1.0, 0.2) -- (1.2, -0.1) -- (1.4, 0.4) -- (1.6, 0.0) 
            -- (1.8, 0.6) -- (2.0, 0.1) -- (2.2, 0.4) -- (2.4, -0.2) 
            -- (2.6, 0.2) -- (2.8, -0.3) -- (3.0, 0.3) -- (3.2, -0.1) 
            -- (3.4, 0.4) -- (3.6, 0.0) -- (3.8, -0.4) -- (4.0, -0.2);
        \fill[blue!70!black] (4, -0.2) circle (1pt);

        % 8. 跨度标注
        \draw[decorate, decoration={brace, amplitude=5pt}] 
            (-0.1, -1.5) -- (-0.1, 0.5) node[midway, left=6pt] {$\Delta$};

        \draw[decorate, decoration={brace, amplitude=5pt, mirror}] 
            (4.1, -0.7) -- (4.1, 1.3) node[midway, right=6pt] {$\Delta$};

        % 标注起始点 r0
        \node[left] at (0, -1.5) {$r_0$};

    \end{tikzpicture}
    \caption{\bilingual{Illustration of the small deviation lower bound within the tilted tube \texorpdfstring{$[r_L, r_L + v\Delta]$}{[r_l, r_l + v Delta]} with end constraints}{带有末端限制 \texorpdfstring{$[r_L, r_L + v\Delta]$}{[r_l, r_l + v Delta]} 的倾斜管道小偏差下界示意图}}
\end{figure}

\bilingual{
    We prove the first inequality \eqref{eq:4.3}, where the event is denoted by $E$. The proof for \eqref{eq:4.3:2} follows from a symmetric construction. To avoid the use of exponential moments, we divide the evolution into two stages, namely a transition stage, whose duration is of order $\Delta^2$, so that the random walk $Y$ can enter the interior of the tube between $\frac{1}{3}$ and $\frac{2}{3}$, and a drifting stage using Mogul'ski\u i's theorem on a sequence of horizontal sub-tubes like a staircase, with each segment lasting of order $\Delta^{2+\eta}$, where $\eta<\varepsilon\wedge 1$, and the last segment is extended to the final time $L$, with duration of order $2\Delta^{2+\eta}$.
}{
    这里证明第一个不等式 \eqref{eq:4.3},记其中的事件为$E$. ~\eqref{eq:4.3:2} 的证明可通过对称构造得出.为了避开指数矩,下面将演化分为两个阶段,即一个过渡阶段,时长的数量级为$\Delta^2$,~ 使得随机游走$Y$能进入到管道内部$\frac 13$至$\frac 23$处,以及在随后阶梯状的一系列水平子管道上分别应用 Mogul'ski\u i 定理的漂移阶段,每一段时长的数量级为$\Delta^{2+\eta}$,~ 其中$\eta<\varepsilon\wedge 1$,~ 最后一段拉长到终点时刻$L$,~ 其数量级为$2\Delta^{2+\eta}$.
}

\begin{proof}\label{proof:lemma:4.3}
    \bilingual{
        Fix $d''\in(\frac{\pi^2}{2}\sigma^2, d')$.
    }{
        固定 $d''\in(\frac{\pi^2}{2}\sigma^2, d')$.
    }

    \bilingual{
        First, we consider the transition from the starting point $0$ to the interior of the tube. Let $t_1 := \lfloor \Delta^2 \rfloor$. We define the transition event as $E_0$, i.e.,
    }{
        首先,考虑从起点 $0$ 进入管道内部的过渡.令 $t_1 := \lfloor \Delta^2 \rfloor$. ~将该过渡事件定义为 $E_0$,~ 即
    }
    \begin{equation*}
        E_0 := \left\{ r_j \le Y_j \le s_j, \, \forall j \in [\![1, t_1]\!],\, Y_{t_1} \in W_1 \right\},
    \end{equation*}
    \bilingual{
        where $W_1 := [ r_{t_1} + \frac{1}{3}\Delta, r_{t_1} + \frac{2}{3}\Delta ]$ is the central window at time $t_1$. By the local limit theorem for random walks in a strip, the probability that the walk stays inside the moving boundaries $r_j\le Y_j\le s_j$ for $j\le t_1$ and ends in the central window $W_1$ is of order $O(1/\sqrt{t_1})=O(\frac 1\Delta)$. Indeed, after scaling, the process converges weakly to a Brownian motion conditioned to stay within $[0,1]$ (starting from $0$), whose terminal distribution has a strictly positive density on $(\frac 13, \frac 23)$. Consequently, 
    }{
        其中 $W_1 := [ r_{t_1} + \frac{1}{3}\Delta, r_{t_1} + \frac{2}{3}\Delta ]$ 是时刻 $t_1$ 的中心窗口.根据条带内随机游走的局部极限定理,游走在 $j\le t_1$ 的移动边界 $r_j\le Y_j\le s_j$ 内保持并且以 $W_1$ 结尾的概率的数量级为 $O(1/\sqrt{t_1})=O(\frac 1\Delta)$. ~事实上,经过缩放后,该过程弱收敛到一个条件布朗运动,该布朗运动被限制在 $[0,1]$ 内 (从 $0$ 开始),其终点分布在 $(\frac 13, \frac 23)$ 上具有严格正密度.因此,
    }
    \begin{equation*}
        \mathbb{P}(E_0) \ge \frac{C}{\Delta} \ge \exp\left(-\frac{d'-d''}{\Delta^2}L\right),
    \end{equation*}
    \bilingual{
        where $C > 0$, and it holds for some $\Delta_1=\Delta_1(d')>0$ and all $\Delta\ge\Delta_1$.
    }{
        其中$C > 0$,~ 对某个$\Delta_1=\Delta_1(d')>0$和所有$\Delta\ge\Delta_1$成立.
    }

    \bilingual{
        Next, we partition the remaining time $L - t_1$. Fix $\eta \in (0, \varepsilon \wedge 1)$ and let the block length be $\ell = \ell(\Delta) := \lfloor \Delta^{2+\eta} \rfloor$, so that we can get $q := \lfloor \frac{L - t_1}{\ell} \rfloor$ blocks. Since $\eta < \varepsilon$, for some $\Delta_2=\Delta_2(\varepsilon,d')>\Delta_1$ and all $\Delta\ge\Delta_2$, we have $q \ge 1$. Let $t_k := t_1 + (k-1)\ell$, $k \in [\![2, q]\!]$, and $t_{q+1} := L$, and let the $k$-th block be $I_k := [\![t_k, t_{k+1}]\!]$, $k\in[\![1, q]\!]$. Note that the last block has length $\ell \le |I_q| = L - t_q < 2\ell$. For each block $I_k$, define the horizontal sub-tube $H_k := [r_{t_{k+1}}, s_{t_k}]$. This tube is well-defined, since $\eta < 1$, so that for some $\Delta_3=\Delta_3(c,\varepsilon)>\Delta_2$ and all $\Delta \ge \Delta_3$, we have $8\frac{c}{\Delta^2}\ell < \Delta$. Hence the tube width is
    }{
        其次,划分剩余的 $L - t_1$. ~固定 $\eta \in (0, \varepsilon \wedge 1)$ 并令分块长度为 $\ell = \ell(\Delta) := \lfloor \Delta^{2+\eta} \rfloor$,~ 从而得到 $q := \lfloor \frac{L - t_1}{\ell} \rfloor$ 个区间. 由于$\eta<\varepsilon$, ~对某个$\Delta_2=\Delta_2(\varepsilon,d')>\Delta_1$和所有$\Delta\ge\Delta_2$成立$q\ge 1$. ~记时刻 $t_k := t_1 + (k-1)\ell$,~ $k \in [\![2, q]\!]$,~ 且 $t_{q+1} := L$,~ 并记第 $k$ 个区间为 $I_k := [\![t_k, t_{k+1}]\!]$,~ $k\in[\![1, q]\!]$. ~注意最后一个区间长度有$\ell\le|I_q|=L-t_q< 2\ell$. ~在每个区间 $I_k$ 内,定义水平子管道 $H_k := [r_{t_{k+1}}, s_{t_k}]$. ~该管道是良定义的,此因$\eta<1$,~ 从而对某个$\Delta_3=\Delta_3(c,\varepsilon)>\Delta_2$和所有 $\Delta \ge \Delta_3$ 成立 $8\frac{c}{\Delta^2}\ell < \Delta$. ~因而管道宽度为
    }
    \begin{equation*}
        \Delta(k) := s_{t_k}-r_{t_{k+1}}
        =-\frac{c}{\Delta^2}(t_{k+1}-t_k) + \Delta
        \ge -\frac{c}{\Delta^2} 2\ell + \Delta > 0,
        \quad
        k\in[\![1, q]\!].
    \end{equation*}
    \bilingual{
        And $H_k\cap H_{k-1}\neq\emptyset$ for all $k\in[\![2, q]\!]$, which is because
    }{
        并且$H_k\cap H_{k-1}\neq\emptyset$对所有$k\in[\![2, q]\!]$成立,此因
    }
    \begin{equation*}
        s_{t_{k-1}}-r_{t_{k+1}}
        =-\frac{c}{\Delta^2}(t_{k+1}-t_{k-1}) + \Delta
        \ge -\frac{c}{\Delta^2} 3\ell + \Delta > 0.
    \end{equation*}
    \bilingual{
        Additionally, at time $t_k$, $k\in[\![2, q]\!]$, define the target window for the interval $I_{k-1}$ as $W_k := [r_{t_{k+1}} + \frac{1}{3}\Delta(k), r_{t_{k+1}} + \frac{2}{3}\Delta(k)]$ and define the final window as $W_{q+1} := [r_L, r_L + v\Delta]$. Note that for $\Delta \ge \Delta_3$, we have $W_1 \subseteq H_1$ and $W_k \subseteq H_k \cap H_{k-1}$ for all $k \in [\![2, q]\!]$, which is because
    }{
        另外在 $t_k$,~ $k\in[\![2, q]\!]$ 处定义区间$I_{k-1}$的目标窗口为 $W_k:=[r_{t_{k+1}} + \frac{1}{3}\Delta(k), r_{t_{k+1}} + \frac{2}{3}\Delta(k)]$且定义最终窗口为 $W_{q+1} := [r_L, r_L + v\Delta]$. ~注意到, 在$\Delta\ge\Delta_3$时,~ $W_1\subseteq H_1$,~ 并且$W_k\subseteq H_k\cap H_{k-1}$ 对所有 $k\in[\![2, q]\!]$ 成立,此因
    }
    \begin{equation*}
        8\frac{c}{\Delta^2}\ell < \Delta
        \implies
        r_{t_{k+1}}
        <
        r_{t_k} + \frac 13\Delta(k)
        <
        r_{t_k} + \frac 23\Delta(k)
        <
        s_{t_{k-1}},
        \quad
        k\in[\![1, q]\!].
    \end{equation*}

    \bilingual{
        For each $k \in [\![1, q]\!]$, let $E_k$ be the event that the walk stays in $H_k$ during $I_k$ and ends in $W_{k+1}$, i.e.,
    }{
        对于每个 $k \in [\![1, q]\!]$,~ 记 $E_k$ 为随机游走 $Y$ 在 $I_k$ 期间保持在 $H_k$ 内且结束于 $W_{k+1}$ 的事件,即
    }
    \begin{equation*}
        E_k := \left\{ r_{t_{k+1}} \le Y_j \le s_{t_k}, \, \forall j \in I_k,\, Y_{t_{k+1}} \in W_{k+1} \right\}.
    \end{equation*}
    \bilingual{
        According to the Markov property, we have
    }{
        根据 Markov 性质,成立
    }
    \begin{align*}
        \mathbb P\left(E_k \middle| \mathcal F_{t_k}\right)
        &=
        \left.\mathbb P\left( \begin{aligned}  - y + r_{t_{k+1}} \le Y_j - y \le  - y + r_{t_{k+1}} + \Delta(k), \, \forall j \in I_k, \\ Y_{t_{k+1}} - y\in -y+W_{k+1} \end{aligned} \right)\right|_{y=Y_{t_k}}
        \\ &=:
        \left.\mathbb P\left(\; \begin{aligned} r'_0 \le \tilde Y_j \le r'_0 + \Delta(k), \, \forall j \in [\![1, |I_k|]\!], \\ \tilde Y_{|I_k|}\in [r'_0 + v_1\Delta(k), r'_0 + v_2\Delta(k)] \end{aligned} \;\right)\right|_{r'_0=-Y_{t_k}+r_{t_{k+1}}},
    \end{align*}
    \bilingual{
        where for any starting position $Y_{t_k} = y \in W_k$, let $\tilde Y_j := Y_{t_k+j} - y$, $j\in[\![1, |I_k|]\!]$ be the random walk shifted to start from $0$ on the interval $I_k$, and let $r'_0 = r'_0(k) := -y + r_{t_{k+1}}$ be the lower bound of the initial position of this random walk, which satisfies $r'_0 \in [-\frac{2}{3}\Delta(k), -\frac{1}{3}\Delta(k)]$. Additionally, let
    }{
        其中,对任意起始位置 $Y_{t_k} = y \in W_k$,~ 记 $\tilde Y_j:=Y_{t_k+j}-y$,~ $j\in[\![1, |I_k|]\!]$ 是区间 $I_k$ 上平移到从 $0$ 开始的随机游走,并记 $r'_0=r'_0(k):=-y+r_{t_{k+1}}$ 是该随机游走的初始位置下界,成立$r'_0\in[-\frac 23\Delta(k),-\frac 13\Delta(k)]$. ~另外,记
    }
    \begin{align*}
        v_1=v_1(\Delta(k)) &:= \begin{cases}
            \frac 1{\Delta(k)}\left(\frac 23\frac c{\Delta^2}|I_{k+1}|+\frac 13\Delta\right), & k\in[\![1, q-1]\!], \\
            0, & k=q
        \end{cases} 
        =O_{\Delta(k)}(1),\\
        v_2=v_2(\Delta(k)) &:= \begin{cases}
            \frac 1{\Delta(k)}\left(\frac 13\frac c{\Delta^2}|I_{k+1}|+\frac 23\Delta\right), & k\in[\![1, q-1]\!], \\
            \frac{v\Delta}{\Delta(k)}, & k=q
        \end{cases}
        =O_{\Delta(k)}(1).
    \end{align*}
    \bilingual{
        Note that, when $\Delta \ge \Delta_3$,
    }{
        注意到,在$\Delta\ge\Delta_3$时,
    }
    $
        8\frac{c}{\Delta^2}\ell < \Delta
        \implies
        v_1(\Delta(k))
        <1.
    $
    \bilingual{
        Apply \nameref{theorem:mogulskii} (see \Cref{lemma:mogulskii-cal} (iii), with $\Delta(k)$ in lieu of $\Delta$, $|I_k|$ in lieu of $L_\Delta$ and $\tilde Y$ in lieu of $Y$), since
    }{
        应用\nameref{theorem:mogulskii}(参见\Cref{lemma:mogulskii-cal} (iii),参数设置为$\Delta(k)$代替$\Delta$,~ $|I_k|$代替$L_\Delta$, ~$\tilde Y$代替$Y$),由于
    }
    \begin{equation*}
        \lim_{\Delta(k)\to\infty} \frac\Delta{\Delta(k)} = 1,
        \quad
        \lim_{\Delta(k)\to\infty} |I_k| = \infty,
        \quad
        \lim_{\Delta(k)\to\infty} \frac{|I_k|}{\Delta(k)^2} = \infty,
    \end{equation*}
    \bilingual{
        and we get that for all $a\in(0,1)$,
    }{
        得到对所有$a\in(0,1)$, 
    }
    \begin{equation}
        % \quad&
        % \lim_{\Delta\to\infty}\frac{\Delta^2}{|I_k|}
        % \log\mathbb P\left(
        %     \begin{aligned}
        %         -a\Delta(k)\le \tilde Y_j\le(1-a)\Delta(k),\,
        %         \forall j\in[\![1,|I_k|]\!],\\
        %         (v_1-a)\Delta(k)\le \tilde Y_{|I_k|}\le(v_2-a)\Delta(k)
        %     \end{aligned}
        % \right)
        % \nonumber \\
        % &
        % =
        \lim_{\Delta(k)\to\infty}
        % \frac{\Delta^2}{\Delta(k)^2}
        \frac{\Delta(k)^2}{|I_k|}
        \log\mathbb P\left(
            \begin{aligned}
                -a\Delta(k)\le \tilde Y_j\le(1-a)\Delta(k),\,
                \forall j\in[\![1,|I_k|]\!],\\
                (v_1-a)\Delta(k)\le \tilde Y_{|I_k|}\le(v_2-a)\Delta(k)
            \end{aligned}
        \right)
        % \nonumber \\ &
        =
        \frac{\pi^2}2\sigma^2.
        \label{eq:lemma:4.3:1}
    \end{equation}
    \bilingual{
        Following the reasoning as in the \nameref{proof:lemma:4.1}, we can deduce that there exists $\Delta_0=\Delta_0(v,d',\varepsilon,c)>\Delta_3$ such that for all $\Delta\ge\Delta_0$, all $k\in[\![1,q]\!]$ and all $r'_0\in[-\frac 23\Delta(k),-\frac 13\Delta(k)]$,
    }{
        参照\nameref{proof:lemma:4.1}的推理,可推得存在 $\Delta_0=\Delta_0(v,d',\varepsilon,c)>\Delta_3$,~ 使得对所有 $\Delta\ge\Delta_0$,~ $k\in[\![1,q]\!]$ 且 $r'_0\in[-\frac 23\Delta(k),-\frac 13\Delta(k)]$, 成立
    }
    \begin{equation*}
        \mathbb P\left(\;
            \begin{aligned}
                r'_0\le \tilde Y_j\le r'_0+\Delta(k),\,
                \forall j\in[\![1,|I_k|]\!],\\
                \tilde Y_{|I_k|}\in[r'_0+v_1\Delta(k),r'_0+v_2\Delta(k)]
            \end{aligned}
        \;\right)
        \ge
        \exp\left(
            -\frac{d''}{\Delta^2}|I_k|
        \right).
    \end{equation*}
    \bilingual{
        Indeed, fix $m\ge 1$ so that $d''{(1-\frac 1{3m})}^2>\frac{\pi^2}2\sigma^2$. Apply~\eqref{eq:lemma:4.3:1} to $(1-\frac 1{3m})\Delta(k)$ in lieu of $\Delta(k)$ and
    }{
        事实上,固定 $m\ge 1$ 使得 $d''{(1-\frac 1{3m})}^2>\frac{\pi^2}2\sigma^2$. ~将\eqref{eq:lemma:4.3:1} 应用于 $(1-\frac 1{3m})\Delta(k)$ 替代 $\Delta(k)$ 和
    }
    \begin{equation*}
        a=\frac{\frac 23-\frac{l+1}{3m}}{1-\frac 1{3m}}\in(0,1),\quad
        l\in[\![0,m-1]\!],
    \end{equation*}
    \bilingual{
        and we can get, for all $l\in[\![0,m-1]\!]$,
    }{
        可得到,对所有 $l\in[\![0,m-1]\!]$,~
    }
    \begin{align*}
        &\quad
        \lim_{\Delta(k)\to\infty}
        \textstyle\frac 
        {
            {\left(1-\frac 1{3m}\right)}^2
            \Delta(k)^2
        }{|I_k|}
        \log\mathbb P\left(
            \begin{aligned}
                \textstyle-1+\frac{\frac{l+1}m+1}3
                \le\frac{\tilde Y_j}{\Delta(k)}
                \le\frac{\frac lm+1}3,\,
                \forall j,\\
                \textstyle\frac{\frac lm+1}3-1+v_1\le
                \frac{\tilde Y_{|I_k|}}{\Delta(k)},\\
                \textstyle\frac{\tilde Y_{|I_k|}}{\Delta(k)}
                \le\frac{\frac lm+1}3-1+v_2
            \end{aligned}
        \right)
        \\
        &=
        \lim_{\Delta\to\infty}
        \textstyle\frac{\Delta(k)^2}{\Delta^2}
        \frac 
        {
            {\left(1-\frac 1{3m}\right)}^2
            \Delta^2
        }{|I_k|}
        \log\mathbb P\left(
            \begin{aligned}
                \textstyle-1+\frac{\frac{l+1}m+1}3
                \le\frac{\tilde Y_j}{\Delta(k)}
                \le\frac{\frac lm+1}3,\,
                \forall j,\\
                \textstyle\frac{\frac lm+1}3-1+v_1\le
                \frac{\tilde Y_{|I_k|}}{\Delta(k)},\\
                \textstyle\frac{\tilde Y_{|I_k|}}{\Delta(k)}
                \le\frac{\frac lm+1}3-1+v_2
            \end{aligned}
        \right)
        \\ &
        =
        -\frac{\pi^2}{2}\sigma^2
        >
        -d''{\left(1-\frac 1{3m}\right)}^2.
    \end{align*}
     \bilingual{
        Then we find $\Delta_0>0$ such that for all $\Delta\ge\Delta_0$ and $l\in[\![0,m-1]\!]$,
    }{
        于是可得 $\Delta_0>\Delta_3$,~ 使得对所有 $\Delta\ge\Delta_0$ 且 $l\in[\![0,m-1]\!]$,
    }
    \begin{equation*}
        \mathbb P\left(
            \begin{aligned}
                \textstyle-1+\frac{\frac{l+1}m+1}3
                \le\frac{\tilde Y_j}{\Delta(k)}
                \le\frac{\frac lm+1}3,\,
                \forall j,\\
                \textstyle\frac{\frac lm+1}3-1+v_1\le
                \frac{\tilde Y_{|I_k|}}{\Delta(k)},\\
                \textstyle\frac{\tilde Y_{|I_k|}}{\Delta(k)}
                \le\frac{\frac lm+1}3-1+v_2
            \end{aligned}
        \right)
        \ge
        \exp\left(
            -\frac{d''}{\Delta^2}|I_k|
        \right).
    \end{equation*}
    \bilingual{
        Now for $r'_0\in[-\frac 23\Delta(k),-\frac 13\Delta(k)]$, we can find some $l\in[\![0,m-1]\!]$ such that $3(\frac{r'_0}{\Delta(k)}+1)-1\in[\frac lm,\frac{l+1}m]$, and we get by monotonicity of probability that
    }{
        现在对 $r'_0\in[-\frac 23\Delta(k),-\frac 13\Delta(k)]$,~ 存在 $l\in[\![0,m-1]\!]$ 使得 $3(\frac{r'_0}{\Delta(k)}+1)-1\in[\frac lm,\frac{l+1}m]$,~ 根据概率的单调性,
    }
    \begin{equation*}
        \mathbb P\left(\;
            \begin{aligned}
                r'_0\le \tilde Y_j\le r'_0+\Delta(k),\,
                \forall j,\\
                r'_0+v_1\Delta(k)\le \tilde Y_{|I_k|},\\
                \tilde Y_{|I_k|}\le r'_0+v_2\Delta(k)
            \end{aligned}
        \;\right)
        \ge
        \mathbb P\left(
            \begin{aligned}
                \textstyle-1+\frac{\frac{l+1}m+1}3
                \le\frac{\tilde Y_j}{\Delta(k)}
                \le\frac{\frac lm+1}3,\,
                \forall j,\\
                \textstyle\frac{\frac lm+1}3-1+v_1\le
                \frac{\tilde Y_{|I_k|}}{\Delta(k)},\\
                \textstyle\frac{\tilde Y_{|I_k|}}{\Delta(k)}
                \le\frac{\frac lm+1}3-1+v_2
            \end{aligned}
        \right)
        \ge
        e^{
            -\frac{d''}{\Delta^2}|I_k|
        },
    \end{equation*}
    \bilingual{
        Finally, using the Markov property in the reverse order, the probability of the event $E$ in \eqref{eq:4.3} is bounded below by the product of probabilities, which yields, for $\Delta\ge\Delta_0$,
    }{
        最后,倒序利用 Markov 性质,~ \eqref{eq:4.3} 中事件 $E$ 的概率被各分块概率之积所下界控制,从而得到,在 $\Delta\ge\Delta_0$ 时，
    }
    \begin{align*}
        \mathbb{P}(E) 
        &\ge
        \mathbb E\left[
            \mathbf 1_{E_0} \prod_{k=1}^{q} \mathbf 1_{E_k}
        \right]
        =
        \mathbb E\left[
            \mathbf 1_{E_0} \prod_{k=1}^{q-1} \mathbf 1_{E_k}
            \mathbb P\left(E_q \middle| \mathcal F_{t_q}\right)
        \right]
        \\
        &
        \ge \mathbb{P}(E_0) \prod_{k=1}^{q} e^{-\frac{d''}{\Delta^2}|I_k|}
        \ge e^{-\frac{d' - d''}{\Delta^2}L} e^{-\frac{d''}{\Delta^2}(L-t_1)}
        \ge e^{-\frac{d'}{\Delta^2}L},
    \end{align*}
    \bilingual{
        which completes the proof of \eqref{eq:4.3}.
    }{
        此即证得 \eqref{eq:4.3}.
    }
    \qedhere
\end{proof}